%% IEEE TAES manuscript wrapper based on the TAES template.
\PassOptionsToPackage{hypertexnames=false}{hyperref}
\documentclass{IEEEtaes}

\usepackage{amsthm}
\usepackage{amsmath}
\usepackage{amssymb}
\usepackage{amsfonts}
\usepackage{array}
\usepackage{booktabs}
\usepackage{float}
\usepackage{graphicx}
\usepackage{tabularx}
\usepackage{textcomp}
\usepackage[urlcolor=blue,linkcolor=blue,citecolor=blue]{hyperref}
\usepackage{algorithm}
\usepackage{algpseudocode}
\usepackage{cite}

\makeatletter
\setlength{\@fptop}{0pt}
\setlength{\@fpsep}{12pt}
\setlength{\@dblfptop}{0pt}
\setlength{\@dblfpsep}{12pt}
\makeatother

\newcommand{\R}{\mathbb{R}}
\newcommand{\dt}{\Delta t}
\newcommand{\norm}[1]{\left\lVert #1 \right\rVert}
\newcommand{\los}{\mathrm{LOS}}
\newcommand{\fov}{\mathrm{FOV}}
\newcommand{\dv}{\Delta v}
\newcommand{\unit}[2]{#1~#2}

\newtheorem{theorem}{Theorem}
\newtheorem{proposition}{Proposition}
\theoremstyle{remark}
\newtheorem{remark}{Remark}
\theoremstyle{plain}

\begin{document}

\title{Viability-Preserving Rollout Approximate Dynamic Programming for Dynamics-Aware Orbital Inspection}

\author{Rugang Tang}
\affil{Northwestern Polytechnical University, Xi'an, China\\ Hong Kong Polytechnic University, Hong Kong, China}

\author{Xin Ning}
\affil{Northwestern Polytechnical University, Xi'an, China}

\author{Chih-Yung Wen}
\affil{Hong Kong Polytechnic University, Hong Kong, China}


\receiveddate{}


\corresp{{\itshape Corresponding author: Rugang Tang}.}

\supplementary{Color versions of one or more of the figures in this article are available online at \href{https://ieeexplore.ieee.org}{https://ieeexplore.ieee.org}.}
\markboth{OrbInspect}{Rollout ADP for Orbital Inspection}

\maketitle

\begin{abstract}
Autonomous orbital inspection couples a history-dependent surface-coverage objective with dynamically reachable and auditable motion. We formulate the problem on a directed candidate graph whose observation nodes store camera-valid target masks and whose safe observable orbital arc (SOOA) edges store finite-duration Hill--Clohessy--Wiltshire transfers and common audit records. The main contribution is a viability-preserving rollout approximate dynamic programming (ADP) policy. At each decision, finite-depth Bellman lookahead evaluates every shield-admissible action and appends a deterministic adaptive safe-greedy completion as an approximate terminal value; branches without a goal-reaching completion receive infinite value. On the deterministic graph, this construction preserves a feasible completion and cannot cost more than its feasible base policy, without requiring a learned critic. A frozen development study compared fitted-critic, heuristic, fixed-rollout, and adaptive-rollout structures and selected depth three before held-out evaluation. On 30 test scenarios generated conditional on incumbent feasibility, the proposed policy and the strongest local-search baseline both reached the 80\% inspectable-coverage goal in every case. Rollout ADP reduced mean $\Delta v$ from 12.946 to \unit{11.239}{m/s} (13.18\%); the paired difference was \unit{-1.707}{m/s} with a 95\% bootstrap interval $[-2.199,-1.265]$, 28 wins, two numerical ties, and no losses ($p=7.45\times10^{-9}$, exact sign test). The $\Delta v$ reduction persisted in 20 out-of-distribution scenarios (8.37\%). Visibility and clearance queries use all 247,525 ISS triangles. The evidence is deterministic and simulator-only and does not establish continuous-time or closed-loop flight safety.
\end{abstract}

\begin{IEEEkeywords}
spacecraft inspection, approximate dynamic programming, rollout, safe observable orbital arc, relative orbital dynamics, coverage planning, Hill--Clohessy--Wiltshire equations
\end{IEEEkeywords}

\section{Introduction}
\label{sec:introduction}

Inspection of large orbital structures is an enabling task for on-orbit servicing and proximity operations. Their external surfaces may contain configuration changes, impact damage, degraded thermal protection, or docking-interface anomalies that fixed onboard cameras cannot observe from a useful geometry \cite{faghihi2022kinodynamic,faghihi2023multiple,apodaca2025trajectory}. An autonomous inspector must therefore decide both where to look and how to move. These decisions are tightly coupled in orbit because a camera pose with high geometric coverage can be unreachable, control-intensive, or unsafe under relative motion.

Existing methods address important parts of this coupling at different resolutions. Robotic next-best-view methods treat sensor range, occlusion, incidence angle, and marginal coverage \cite{scott2003view}; kinodynamic coverage planners integrate observation and differential constraints in continuous state space \cite{faghihi2022kinodynamic,faghihi2023multiple}; and relative-orbit planners exploit natural orbital motion \cite{hou2024relative,deka2023astrodynamics}. The present problem additionally requires a finite, auditable action representation for repeated sequence decisions under a hard coverage goal. A view-first pipeline that checks transfers only after selection can overestimate useful coverage: individually attractive views may be connected only by motions that violate input or keep-out limits.

The deeper issue is the definition of a selectable inspection action. A static viewpoint is incomplete because it contains no evidence that the spacecraft can reach that pose under its relative dynamics. A dynamically natural orbit segment is also incomplete when it does not reveal new weighted surface area. For orbital inspection, the action itself should include the transfer, the stabilized camera observation, and the safety evidence used to admit that action. This requirement motivates a finite representation that binds observability and relative-motion feasibility before the coverage sequence is optimized.

We introduce the safe observable orbital arc (SOOA), a source-dependent finite-duration HCW transfer from the current graph node to a stabilized candidate observation. Each SOOA references its destination's camera-valid target set and stores control effort, minimum sampled clearance, terminal error, input and speed status, and optional passive-drift audit. Geometric viewpoints define graph nodes, but they do not contribute coverage until the corresponding directed SOOA has passed the enabled feasibility tests. Figure~\ref{fig:sooa-framework} summarizes the progression from mesh targets and camera visibility to a directed graph of admissible inspection actions. The graph objective then trades newly covered weighted area against the cost of each reachable arc.

The finite representation makes sequential optimization and verification explicit. A node identifier alone is not Markov because the value of a view depends on what has already been covered and how much budget remains. We therefore augment the state with covered- and selected-set masks and develop viability-preserving rollout ADP. A shield-admissible safe-greedy policy supplies an approximate cost-to-go; finite-depth Bellman expansion improves the first action, the state is updated, and the procedure repeats until the coverage goal is reached \cite{bertsekas2020rollout}. Exact dynamic programming remains available only for reduced graphs small enough to exhaust.

The contributions are:
\begin{enumerate}
  \item a node--edge representation in which candidate observations store camera-valid target sets and directed SOOAs store source-dependent HCW motion and sampled audit evidence;
  \item a full-mesh bounding-volume hierarchy used consistently for camera line of sight, candidate clearance, sampled trajectory clearance, and swept surface-crossing rejection;
  \item a viability-preserving rollout ADP algorithm that combines finite-depth Bellman improvement, an adaptive safe-greedy terminal policy, and infinite value for states without a goal-reaching completion;
  \item proofs of Markov sufficiency, compositional preservation of stored audits, rollout viability and base-policy no-degradation, and exact optimality on exhausted reduced graphs; and
  \item a split-controlled development and held-out study that reports the unsuccessful fitted critic as a negative ablation, freezes rollout depth on validation, and reports paired effect sizes, confidence intervals, exact sign tests, distribution-shift results, and exact non-timing reproducibility.
\end{enumerate}

\begin{figure*}[t]
  \centering
  \includegraphics[width=\textwidth]{figures/orbinspect_sooa_framework_full_mesh.pdf}
  \caption{OrbInspect formulation. (a) Candidate observations are generated around an ISS-scale structure and outside a mesh-derived safety margin. (b) Range, field of view, incidence, full-mesh line of sight, and the prescribed observation condition determine each node's visible-target set. (c) A directed SOOA edge binds a source-dependent HCW transfer and its sampled and swept-mesh audits to a destination observation. (d) Rollout ADP expands only shield-admissible choices and rejects coverage-rich actions without a viable completion.}
  \label{fig:sooa-framework}
\end{figure*}

\section{Related Work}
\label{sec:related}

Relative-motion models provide the dynamic foundation for orbital inspection. The Clohessy--Wiltshire equations remain a common baseline near a circular chief orbit because they are linear, interpretable, and inexpensive to propagate \cite{clohessy1960terminal,wertz2011space}. Higher-fidelity flight design must account for nonlinear gravity, perturbations, navigation error, and actuator dynamics, but HCW propagation is useful for isolating how view ordering changes control effort and safety under a fixed model. Visibility-constrained orbit design has also been studied directly. Hou et al. analyze relative orbits for visual inspection of the China Space Station \cite{hou2024relative}, while Zhou et al. combine near-natural relative ellipses with formation control around a rotating station model \cite{zhou2026formation}. These approaches provide structured orbital motion and stronger closed-loop control analysis than the present planner, but their primary object is an orbit or formation rather than a sequence of surface-coverage actions.

Robotic inspection provides the complementary view-selection perspective. Classical next-best-view work models sensor range, incidence, occlusion, reconstruction quality, and motion cost \cite{scott2003view}. Sampling-based kinodynamic planners such as RRT and asymptotically optimal variants incorporate differential constraints and collision avoidance in continuous state space \cite{lavalle1998rapidly,karaman2011sampling}. Spacecraft-specific extensions combine full-coverage inspection with kinodynamic planning for one or several inspectors \cite{faghihi2022kinodynamic,faghihi2023multiple}, and astrodynamics-informed sampling exploits natural relative motion during transfer search \cite{deka2023astrodynamics}. These methods are valuable when continuous path generation, coordination, or obstacle avoidance dominates. OrbInspect instead treats each audited transfer-observation pair as a finite decision action and carries the accumulated weighted coverage explicitly in the decision state.

Approximate dynamic programming and rollout address sequential decisions when exact Bellman recursion is too large \cite{bertsekas2020rollout,bertsekas1996neuro}. Rollout is an ADP method even without a learned critic: the cost-to-go of a computationally inexpensive base policy approximates the exact value, and a truncated Bellman minimization performs online policy improvement. Learning-based policies additionally require a safety argument: constrained policy optimization controls expected constraint costs \cite{achiam2017cpo}, whereas shielding removes actions that violate an independently specified safety model \cite{alshiekh2018shielding}. We combine the latter separation with a viability-aware rollout value. A fitted linear critic was tested as an alternative architecture, but it failed the predeclared success gate and is reported as a negative ablation rather than presented as the proposed method.

Recent spacecraft-inspection studies increasingly couple geometry and motion. Apodaca et al. address inspection trajectory generation and coverage--energy trade-offs for in-orbit structures \cite{apodaca2025trajectory}. Zhou et al. study surface-aware tomographic paths around complex spacecraft geometry \cite{zhou2026tomographic}. Ogri et al. address an execution-level problem involving image-feature tracking, information gain, and dwell under switched-system logic \cite{ogri2025switched}. Together, these studies show that geometric coverage, orbital dynamics, and online sensing cannot be treated as fully independent. They also differ in the level at which the coupling occurs: orbit design, continuous path generation, online camera control, or discrete inspection planning.

OrbInspect is positioned at the discrete-action level. Its contribution is neither a new relative dynamics model nor a flight-certified collision-avoidance controller. Each SOOA is simultaneously a sensor-valid coverage contributor and a dynamically audited finite motion, so an infeasible transfer cannot improve the planner's reported coverage. The algorithmic contribution is the coupling of rollout ADP with a finite-completion viability test on this inspection graph, not a claim that rollout itself is new. Formation-level studies provide complementary lessons: near-natural relative arcs can seed lower-effort SOOA libraries, and a rotating body frame should replace the fixed-LVLH target assumption when station attitude motion is material. All optimality and safety statements remain conditional on the sampled library and stated audits.

\section{System and Observation Model}
\label{sec:system}

The present study addresses offline planning. It receives a known station mesh, an initial relative state, camera parameters, a keep-out model, and actuation limits, and returns a time-indexed translational reference, a camera-boresight stream, selected SOOA records, and summary metrics. A future ROS~2 replay layer may publish these records, but no ROS or Gazebo result is used to support the quantitative claims in this manuscript. Table~\ref{tab:assumptions} states the modeling assumptions and their implications for interpreting the results.

\begin{table*}[t]
  \centering
  \caption{Modeling assumptions and the boundary of the reported evidence.}
  \label{tab:assumptions}
  \renewcommand{\arraystretch}{1.08}
  \setlength{\tabcolsep}{5pt}
  \footnotesize
  \begin{tabularx}{\textwidth}{@{}p{0.18\textwidth}p{0.34\textwidth}X@{}}
    \toprule
    Assumption & Role in the baseline & Consequence for interpretation \\
    \midrule
    Circular chief orbit & Supports linear HCW relative translation in the LVLH frame. & Quantitative results compare planners under the same local linear model, not a perturbed flight ephemeris. \\
    Station fixed in LVLH & Keeps mesh targets and the mesh safety geometry stationary during one planning run. & Rotating appendages or station attitude motion require a body-frame extension. \\
    Known mesh and state & Makes visibility, occlusion, and transfer evaluation deterministic. & Navigation, map, and shape uncertainty are not represented by the present safety margin. \\
    Prescribed camera attitude & Points the boresight toward each candidate observation target and exports the corresponding quaternion stream. & Camera slew time, reaction-wheel limits, and closed-loop attitude error are outside the optimization. \\
    Mesh-based sampled safety & Computes distance to the full ISS triangle mesh at every rollout sample and rejects sampled motion segments that intersect the mesh. & Positive sampled clearance and surface nonintersection are planning audits rather than a robust continuous-time collision certificate. \\
    No illumination or plume model & Isolates view geometry and translational feasibility. & Sun angle, shadow, contamination, and plume-impingement restrictions must be added for flight design. \\
    \bottomrule
  \end{tabularx}
\end{table*}

\subsection{Relative Dynamics}

The inspected structure is fixed in an LVLH frame attached to a nominal circular chief orbit. The axes follow the radial, along-track, and cross-track directions shown in Fig.~\ref{fig:sooa-framework}. The chaser state is $\mathbf{x}=[\mathbf{r}^{\mathsf T},\mathbf{v}^{\mathsf T}]^{\mathsf T}\in\R^6$, where $\mathbf r=[r_x,r_y,r_z]^{\mathsf T}$ and $\mathbf v=[v_x,v_y,v_z]^{\mathsf T}$. With mean motion $n$ and commanded acceleration $\mathbf{u}=[a_x,a_y,a_z]^{\mathsf T}$, the HCW model is
\begin{equation}
  \begin{aligned}
  \dot{\mathbf{x}}&=A_{\mathrm{HCW}}\mathbf{x}+B\mathbf{u},\\
  A_{\mathrm{HCW}}&=\begin{bmatrix}0&I\\K_n&C_n\end{bmatrix},
  &B&=\begin{bmatrix}0\\I\end{bmatrix},\\
  K_n&=\operatorname{diag}(3n^2,0,-n^2),
  &C_n&=\begin{bmatrix}0&2n&0\\-2n&0&0\\0&0&0\end{bmatrix}.
  \end{aligned}
  \label{eq:hcw-continuous}
\end{equation}
For the continuous HCW model, piecewise-constant acceleration over $\dt$ has the exact zero-order-hold map
\begin{equation}
  \begin{aligned}
  \mathbf{x}_{k+1}&=A_d\mathbf{x}_k+B_d\mathbf{u}_k,\\
  A_d&=e^{A_{\mathrm{HCW}}\dt},\\
  B_d&=\int_0^{\dt}e^{A_{\mathrm{HCW}}\tau}B\,d\tau.
  \end{aligned}
  \label{eq:hcw-discrete}
\end{equation}
The reported implementation uses the same piecewise-constant input model but constructs and propagates each discrete boundary map with fixed-step fourth-order Runge--Kutta integration. We denote that numerical affine map by $\mathbf x_{k+1}=\widetilde A_d\mathbf x_k+\widetilde B_d\mathbf u_k$; Eq.~\eqref{eq:hcw-discrete} is the exact continuous-model reference, not a claim that the archived samples were generated by a matrix exponential. The transfer is therefore a finite-duration continuous-thrust maneuver with zero-order-held acceleration, which may jump at sample boundaries, rather than an impulsive maneuver. Each segment is solved rest-to-rest at a candidate observation, and the numerical rollout is accepted only when its terminal position and velocity errors satisfy the stated tolerances.

\subsection{Visibility, Prescribed Observation, and Coverage}

The structure is represented by weighted surface samples $\mathcal T=\{(p_i,n_i,w_i)\}_{i=1}^{N}$ with $w_i>0$. Each weight approximates the surface area represented by its sample, so a densely tessellated component does not dominate solely because it contains more vertices. The camera attitude $q_k$ defines unit boresight $b(q_k)$, and $\mathcal O$ denotes the occluding mesh. Target $i$ is visible when
\begin{align}
  r_{\min}^{\mathrm{cam}}&\le\norm{p_i-\mathbf r_k}\le r_{\max}^{\mathrm{cam}},\nonumber\\
  \arccos\!\frac{(p_i-\mathbf r_k)^{\mathsf T}b(q_k)}
  {\norm{p_i-\mathbf r_k}}&\le\alpha_{\max},\nonumber\\
  \arccos\!\frac{(\mathbf r_k-p_i)^{\mathsf T}n_i}
  {\norm{\mathbf r_k-p_i}\norm{n_i}}&\le\theta_{\max},\nonumber\\
  \los(\mathbf r_k,p_i,\mathcal O)&=1.
  \label{eq:visibility}
\end{align}
The four tests exclude targets that are too near or too far, outside the field of view, observed at an excessive incidence angle, or blocked by station geometry. Line of sight is queried against a deterministic bounding-volume hierarchy containing every ISS mesh triangle; intersections at the target endpoint are excluded, while interior segment intersections block the view. In a closed-loop execution model, dwell could be accumulated as $\tau_{i,k+1}=\tau_{i,k}+\dt$ while target $i$ remains visible and accepted when $\tau_i\ge\tau_{\min}$. The offline study instead treats every candidate as a stabilized terminal observation with its prescribed exposure completed a priori. Selecting node $j$ therefore credits its fixed visibility mask $G_j$, so $m_{k+1}=m_k\lor G_j$. The weighted coverage represented by mask $m_k$ is
\begin{equation}
  C(m_k)=\frac{\sum_i w_i m_{i,k}}{\sum_i w_i}.
  \label{eq:coverage}
\end{equation}

Candidate generation can make a subset of the sampled surface invisible from every candidate observation. For the nonempty set $\mathcal T_{\mathrm{insp}}=\cup_{j=1}^{M}G_j$, we therefore also report inspectable coverage:
\begin{equation}
  C_k^{\mathrm{insp}}=
  \frac{\sum_{i\in\mathcal T_{\mathrm{insp}}}w_i m_{i,k}}
  {\sum_{i\in\mathcal T_{\mathrm{insp}}}w_i}.
  \label{eq:inspectable-coverage}
\end{equation}
The total coverage in Eq.~\eqref{eq:coverage} measures how much of the sampled structure is credited, whereas Eq.~\eqref{eq:inspectable-coverage} conditions on the candidate library. The reported mission duration includes transfer time but no additional dwell or exposure time. Attitude determines the visibility predicate and exported boresight; closed-loop attitude dynamics, illumination, and interruption of a dwell are outside the offline model.

\subsection{Dynamic Feasibility and SOOAs}

Every sampled state must satisfy
\begin{equation}
  \norm{\mathbf u_k}_2\le u_{\max},\qquad
  \norm{\mathbf v_k}_2\le v_{\max},\qquad
  d_{\mathcal S}(\mathbf r_k)\ge 0,
  \label{eq:limits}
\end{equation}
where $d_{\mathcal S}(\mathbf r)=d_{\mathrm{mesh}}(\mathbf r)-d_{\mathrm{safe}}$ is clearance above the required margin and $d_{\mathrm{mesh}}$ is unsigned distance to the closest triangle in the complete ISS mesh; thus $d_{\mathcal S}=0$ is the sampled admissibility boundary. In addition, each consecutive pair of rollout samples is tested for a triangle-surface intersection. This swept test catches a direct mesh crossing even if both endpoints have positive unsigned distance. For a directed rollout from source node $i$ to destination node $j$, define $\Delta v_{ij}=\sum_k\norm{\mathbf u_{ij,k}}_2\dt$, root-mean-square distance $e_{ij}^{\mathrm{rms}}$ from the transfer samples to the destination position, and $d_{ij}^{\min}=\min_k d_{\mathcal S}(\mathbf r_{ij,k})$. The graph stage cost used in the reported study is
\begin{equation}
  \ell_{ij}=\lambda_v\Delta v_{ij}+\lambda_e e_{ij}^{\mathrm{rms}}+\lambda_s\max(0,-d_{ij}^{\min})+\lambda_a.
  \label{eq:edge-cost}
\end{equation}
For every shield-admissible edge, the clearance-violation term is zero. It remains in the shared cost implementation so deliberately unshielded comparison methods can be evaluated with the same expression.

A candidate observation node and a directed SOOA edge are, respectively,
\begin{equation}
  \begin{aligned}
  v_j&=(\bar{\mathbf r}_j,\bar q_j,G_j),\\
  a_{ij}&=\bigl(\mathbf x_{ij}[1{:}K],\mathbf u_{ij}[0{:}K-1],q_{ij}[1{:}K],\\
  &\qquad \Delta v_{ij},d_{ij}^{\min},u_{ij}^{\max},v_{ij}^{\max},e_{ij}^{r},e_{ij}^{v},\rho_{ij}\bigr).
  \end{aligned}
  \label{eq:sooa}
\end{equation}
Here $\bar{\mathbf r}_j$ and $\bar q_j$ are the terminal observation pose and prescribed attitude, $G_j\subseteq\mathcal T$ is its camera-valid target set, and $e_{ij}^{r}$ and $e_{ij}^{v}$ are terminal position and velocity errors. The stored state and attitude arrays contain the $K$ post-step samples; the source state is held by node $v_i$. The initial condition is a distinguished source node $v_0$ with no visibility reward. For passive horizon $T_p=L_p\dt$, an optional audit can propagate zero-input HCW drift from every stored edge state and record $\rho_{ij}=\min_{1\le k\le K,\,0\le l\le L_p}d_{\mathcal S}(\mathbf r_{ij,k}^{\mathrm{drift}}(l))$. The held-out study disables this optional audit and therefore makes no passive-safety claim; it retains input, speed, full-mesh sampled clearance, swept surface-intersection, and terminal checks. Throughout this paper, ``safe'' means admissible under the audits explicitly enabled for the reported experiment. These deterministic checks do not certify robustness between samples or under model and state uncertainty.

\section{Viability-Preserving Rollout ADP}
\label{sec:algorithm}

\subsection{Candidate Graph and Markov State}

Let $\mathcal V=\{v_0,v_1,\ldots,v_M\}$ contain the initial node and the $M$ candidate observation nodes in Eq.~\eqref{eq:sooa}. A directed edge from $v_i$ to $v_j$ is the source-dependent SOOA $a_{ij}$, and it is absent when an enabled audit fails. Thus visibility $G_j$ is node dependent, whereas trajectory, cost, clearance, input, terminal error, and passive margin are edge dependent. This separation is necessary because reaching the same observation from two different source nodes generally produces different HCW motion and audit values.

A node identifier is not a sufficient decision state because the value of a view changes after its targets have been covered. The finite-horizon Markov state is
\begin{equation}
  s_k=(j_k,m_k,b_k,h_k),
  \label{eq:mdp-state}
\end{equation}
where $j_k\in\{0,\ldots,M\}$ is the current node, $m_k\in\{0,1\}^{N}$ is the covered-target mask, $b_k\in\{0,1\}^{M}$ is the selected-observation mask, and $h_k$ is the remaining action budget. Selecting destination $a\in\{1,\ldots,M\}$ traverses edge $a_{j_ka}$ and gives the deterministic transition
\begin{equation}
  f(s_k,a)=\left(a,\;m_k\lor G_a,\;b_k\lor\mathbf e_a,\;h_k-1\right).
  \label{eq:mdp-transition}
\end{equation}
The shield-admissible action set $\mathcal U_s(s)$ is defined below. The constrained graph problem assigns zero terminal cost after reaching $C_{\mathrm{req}}^{\mathrm{insp}}$ and value $+\infty$ when no goal-reaching continuation exists within the remaining budget. Its Bellman recursion is
\begin{equation}
  \begin{aligned}
  V^\star(s)&=
    \begin{cases}
    0, & C^{\mathrm{insp}}(m)\ge C_{\mathrm{req}}^{\mathrm{insp}},\\
    +\infty, & h=0\ \text{or}\ \mathcal U_s(s)=\varnothing,\\
    \min\limits_{a\in\mathcal U_s(s)}
      Q^\star(s,a), & \text{otherwise},
    \end{cases}\\
  Q^\star(s,a)&=\ell_{j a}+V^\star(f(s,a)).
  \end{aligned}
  \label{eq:bellman}
\end{equation}
Here $j$ and $h$ are the current-node and remaining-budget components of $s$. Exact evaluation is exponential in the coverage and selected-node masks. The proposed method approximates this recursion with a short exact prefix and a computationally inexpensive, goal-aware terminal policy.

\subsection{SOOA Construction and Safety Shield}

Area-weighted targets are sampled from the mesh, and camera seeds are generated on offset shells around their outward normals. The binary visibility matrix $V_{ij}=\mathbf 1\{i\in G_j\}$ is computed once from the terminal observation geometry. For a transition to $\bar{\mathbf x}_j=[\bar{\mathbf r}_j^{\mathsf T},\mathbf0^{\mathsf T}]^{\mathsf T}$, the planner solves the minimum-energy piecewise-constant HCW boundary-value problem under the numerical RK4 map
\begin{equation}
  \begin{aligned}
  \min_{\mathbf u_0,\ldots,\mathbf u_{K-1}}\quad&
  \sum_{k=0}^{K-1}\norm{\mathbf u_k}_2^2\\
  \mathrm{s.t.}\quad&
  \mathbf x_{k+1}=\widetilde A_d\mathbf x_k+\widetilde B_d\mathbf u_k,\\
  &\mathbf x_K=\bar{\mathbf x}_j.
  \end{aligned}
  \label{eq:terminal-transfer}
\end{equation}
If the terminal influence matrix $\mathcal B_K=[\widetilde A_d^{K-1}\widetilde B_d\ \cdots\ \widetilde B_d]$ has full row rank, the stacked minimum-norm command from source state $\mathbf x_i^{\mathrm{src}}$ is $\mathbf U_{ij}^\star=\mathcal B_K^{\mathsf T}(\mathcal B_K\mathcal B_K^{\mathsf T})^{-1}(\bar{\mathbf x}_j-\widetilde A_d^K\mathbf x_i^{\mathrm{src}})$. The implementation computes this expression through the $6\times6$ Gram system and then propagates the requested input without clipping. Let $\chi_{ij}=1$ exactly when every stored edge sample satisfies Eq.~\eqref{eq:limits}, $e_{ij}^{r}\le\varepsilon_r$, $e_{ij}^{v}\le\varepsilon_v$, and $\rho_{ij}\ge0$ when the passive audit is enabled. The shield then admits
\begin{equation}
  \begin{split}
  \mathcal U_s(s)=\{a\in\{1,\ldots,M\}:{}& b_a=0,\;\chi_{ja}=1,\\
    &G_a\setminus\operatorname{supp}(m)\neq\varnothing\},
  \end{split}
  \label{eq:safe-action-set}
\end{equation}
where $j$ is the current-node component of $s$. Unlike a reward penalty, the shield makes an inadmissible edge unavailable to incumbent construction, local search, and exact recursion \cite{alshiekh2018shielding}. The guarantee remains conditional on the sampled numerical rollout and the deterministic model used to compute $\chi_{ij}$.

\subsection{Adaptive Base Policy and Rollout Value}

Let $g(s,a)=C^{\mathrm{insp}}(m\lor G_a)-C^{\mathrm{insp}}(m)$ denote the marginal inspectable-coverage gain of action $a$, and let $\epsilon=0.05$ prevent division by a near-zero stage cost. The deterministic adaptive base policy is
\begin{equation}
  \mu(s)\in\arg\max_{a\in\mathcal U_s(s)}
  \left(\frac{g(s,a)}{\max\{\epsilon,\ell_{ja}\}},
  g(s,a),-\ell_{ja}\right),
  \label{eq:base-policy}
\end{equation}
where the tuple is compared lexicographically, followed by deterministic node-identifier tie breaking. The policy recomputes gain after every accepted observation. Starting from state $s$, its completion value is
\begin{equation}
  \widehat V_0(s)=
  \begin{cases}
  \sum_{k=0}^{L-1}\ell_{j_k\mu(s_k)}, & \mu\text{ reaches }C_{\mathrm{req}}^{\mathrm{insp}}\text{ within }h,\\
  +\infty, & \text{otherwise}.
  \end{cases}
  \label{eq:base-value}
\end{equation}
Thus $\widehat V_0$ is not a learned critic. It is a model-based approximate cost-to-go obtained by simulating a complete, shielded base policy.

For rollout depth $d\ge1$, define the truncated Bellman value
\begin{equation}
  \widehat V_d(s)=
  \begin{cases}
    0, & C^{\mathrm{insp}}(m)\ge C_{\mathrm{req}}^{\mathrm{insp}},\\
    +\infty, & h=0,\\
    \min\limits_{a\in\mathcal U_s(s)}\widehat Q_d(s,a), & \text{otherwise},
  \end{cases}
  \label{eq:rollout-value}
\end{equation}
where $\widehat Q_d(s,a)=\ell_{ja}+\widehat V_{d-1}(f(s,a))$.
Only actions with finite completion value participate in the minimization. The rollout policy selects the minimizing first action, applies Eq.~\eqref{eq:mdp-transition}, and rebuilds the depth-$d$ tree at the new state. This repeated policy-improvement step distinguishes the proposed method from following a fixed incumbent suffix. State--depth memoization removes duplicate subproblems within one decision.
Figure~\ref{fig:adp-architecture} shows how the finite-completion test, safety shield, and repeated depth-$d$ improvement interact.

\begin{figure*}[t]
  \centering
  \includegraphics[width=\textwidth]{figures/adp_future/adp_rollout_architecture.pdf}
  \caption{Proposed viability-preserving rollout ADP. The finite candidate-graph state carries current node, covered and selected masks, and remaining budget. The safety shield removes inadmissible SOOAs before search. A depth-$d$ exact prefix is terminated by a complete adaptive safe-greedy simulation; a branch is viable only if this completion reaches the coverage goal. The first minimizing action is executed and the value is recomputed.}
  \label{fig:adp-architecture}
\end{figure*}

\begin{algorithm}[t]
\caption{Viability-preserving rollout ADP}
\label{alg:safe-graph-api}
\begin{algorithmic}[1]
\Require Candidate graph, shield $\mathcal U_s$, goal $C_{\mathrm{req}}^{\mathrm{insp}}$, budget $H$, depth $d$
\State $s\gets(0,\mathbf0,\mathbf0,H)$; $\pi\gets()$
\While{$C^{\mathrm{insp}}(m)<C_{\mathrm{req}}^{\mathrm{insp}}$ and $h>0$}
  \State Enumerate every $a\in\mathcal U_s(s)$
  \For{each admissible $a$}
    \State $\widehat Q_d(s,a)\gets\ell_{ja}+\widehat V_{d-1}(f(s,a))$
  \EndFor
  \State Discard every action with $\widehat Q_d(s,a)=+\infty$
  \If{no finite action remains} \State \Return infeasible \EndIf
  \State $a^\star\gets\arg\min_a \widehat Q_d(s,a)$ with deterministic ties
  \State Append $a^\star$ and its audit record to $\pi$
  \State $s\gets f(s,a^\star)$
\EndWhile
\State \Return $\pi$ if the goal is reached; otherwise return infeasible
\end{algorithmic}
\end{algorithm}

\subsection{Properties and Exact Oracle}
\label{sec:theory}

\begin{proposition}[Monotone submodularity of fixed-mask coverage]
\label{prop:submodular-coverage}
For an observation set $S\subseteq\{1,\ldots,M\}$, let $F(S)=\sum_{i=1}^{N}w_i\mathbf1\{i\in\cup_{j\in S}G_j\}$. If all $w_i\ge0$, then $F$ is normalized, monotone, and submodular.
\end{proposition}

\begin{IEEEproof}
Clearly $F(\varnothing)=0$. If $A\subseteq B$, every target covered by $A$ is covered by $B$, so nonnegative weights give $F(A)\le F(B)$. For $a\notin B$, the marginal gain $F(A\cup\{a\})-F(A)$ is the total weight in $G_a\setminus\cup_{j\in A}G_j$. Since $A\subseteq B$, this set contains $G_a\setminus\cup_{j\in B}G_j$, proving diminishing returns.
\end{IEEEproof}

\begin{remark}[Role of submodularity]
Proposition~\ref{prop:submodular-coverage} justifies marginal-coverage screening at the node level. It does not supply a greedy approximation ratio for the complete inspection tour because the directed SOOA cost and admissibility of the next observation depend on the current node.
\end{remark}

\begin{proposition}[Markov sufficiency]
\label{prop:markov-sufficiency}
For a fixed deterministic candidate graph with fixed node masks, source-dependent edge records, target weights, revisit rule, and horizon, the state in Eq.~\eqref{eq:mdp-state} is Markov for the finite graph decision problem.
\end{proposition}

\begin{IEEEproof}
The current node $j_k$ identifies the source of every possible next edge; $m_k$ determines marginal and terminal coverage; $b_k$ determines which observations remain selectable under the no-revisit rule; and $h_k$ determines the remaining budget. By assumption, visibility masks and edge costs and audits are deterministic records indexed only by the current and destination nodes. Equations~\eqref{eq:mdp-transition}, \eqref{eq:bellman}, and \eqref{eq:safe-action-set} therefore depend on no earlier history once $(j_k,m_k,b_k,h_k)$ is given. The proposition concerns this finite graph abstraction, not the underlying uncertain continuous spacecraft system.
\end{IEEEproof}

\begin{proposition}[Compositional preservation of stored audits]
\label{prop:audit-composition}
Consider a sequence $(a_{j_{k-1}j_k})_{k=1}^{L}$ whose first source is $v_0$, whose idealized rest-to-rest endpoints match exactly, and whose every edge has $\chi_{j_{k-1}j_k}=1$. Then every stored sample of the concatenated sequence satisfies the input, speed, and clearance inequalities, every edge endpoint satisfies its terminal conditions, and every stored passive-drift rollout has nonnegative margin when that audit is enabled.
\end{proposition}

\begin{IEEEproof}
By the definition of $\chi_{ij}$, all stated conditions hold on each accepted edge record. Exact endpoint matching identifies the terminal state of edge $k$ with the source state of edge $k+1$, so the records concatenate without an unmodelled jump. The sample set of the complete trajectory is the finite union of the per-edge sample sets; therefore every member of that union retains the corresponding per-edge audit result.
\end{IEEEproof}

\begin{remark}[Meaning of audit preservation]
The proposition is a finite-record composition result, not a continuous-time barrier certificate. It does not cover inter-sample clearance, endpoint perturbations beyond tolerance, closed-loop tracking error, or robustness to model and state uncertainty.
\end{remark}

\begin{proposition}[Rollout viability]
\label{prop:rollout-viability}
For the deterministic finite graph, if $\widehat V_d(s)<+\infty$, Algorithm~\ref{alg:safe-graph-api} selects a shield-admissible action whose successor has a finite goal-reaching completion within the remaining budget. Replanning preserves this property until the goal is reached.
\end{proposition}

\begin{IEEEproof}
Finiteness in Eq.~\eqref{eq:rollout-value} provides a certificate consisting of shield-admissible prefix actions followed, unless the goal is reached earlier, by a base-policy completion. The selected first action is part of such a certificate, so its successor retains a goal-reaching sequence within the reduced budget. When the depth-$d$ tree is rebuilt, the certificate tail remains admissible; if another exact prefix action is required, the next action of the certified base completion supplies it. Determinism preserves the recorded masks, costs, and audits. Induction over the decreasing action budget therefore proves goal attainment.
\end{IEEEproof}

\begin{proposition}[Base-policy no-degradation]
\label{prop:base-bound}
If the adaptive base policy reaches the goal from $s$, then $\widehat V_d(s)$ is finite and $\widehat V_d(s)\le\widehat V_0(s)$ for every $d\ge1$. For a fixed depth $d$, the realized graph cost of repeated replanning in Algorithm~\ref{alg:safe-graph-api} is no greater than $\widehat V_d(s)$ and therefore no greater than the base-policy cost.
\end{proposition}

\begin{IEEEproof}
The first $d$ base-policy actions, truncated if the goal is reached earlier, followed by the remaining base completion form a feasible certificate in Eq.~\eqref{eq:rollout-value}; hence $\widehat V_d(s)\le\widehat V_0(s)$. Now take a minimizing depth-$d$ certificate at a replanning state and let $s'$ be the successor of its first action $a$. The certificate tail, extended when necessary by the next base-policy action, is feasible in the rebuilt depth-$d$ tree at $s'$. Therefore $\ell_{ja}+\widehat V_d(s')\le\widehat V_d(s)$. Summing this inequality over successive replans telescopes to a realized cost no greater than the initial $\widehat V_d(s)$. The result concerns the graph stage cost; reductions in physical $\Delta v$ are evaluated separately.
\end{IEEEproof}

For reduced-graph certification, exact recursion caches Eq.~\eqref{eq:bellman} for every reachable tuple $(j,m,b,h)$ and enumerates every action in $\mathcal U_s(s)$ rather than a local reversal neighborhood.

\begin{theorem}[Optimality on an exhausted reduced SOOA graph]
\label{thm:reduced-optimality}
For fixed finite nodes, deterministic visibility masks, deterministic audited edge costs, a coverage requirement, and horizon $H$, exhaustive evaluation of Eq.~\eqref{eq:bellman} returns a minimum-cost goal-reaching shield-feasible sequence when one exists. If no such sequence exists, it returns value $+\infty$ and reports infeasibility.
\end{theorem}

\begin{IEEEproof}
Proceed by induction on the remaining budget $h$. At $h=0$, Eq.~\eqref{eq:bellman} returns zero exactly for a goal state and $+\infty$ otherwise, which is correct. Assume the recursion is correct for all states with at most $h-1$ remaining actions. At a state with $h$ actions, every feasible sequence begins with exactly one $a\in\mathcal U_s(s)$ and then a feasible continuation from $f(s,a)$; by the induction hypothesis, the recursion assigns that continuation its minimum cost or $+\infty$ if none exists. Minimizing over every admissible first action therefore returns the minimum total cost, or $+\infty$ when all continuations are infeasible. The claim is limited to the exhausted sampled graph and does not extend to omitted camera poses or continuous-space trajectories.
\end{IEEEproof}

\begin{remark}[Graph certificate versus mission optimality]
Theorem~\ref{thm:reduced-optimality} certifies only the nodes, edges, masks, and audits included in the exhausted graph. Omitted camera poses, alternative transfer durations, and continuous-space trajectories remain outside the certified decision set.
\end{remark}

With maximum branching factor $B$, depth $d$, and at most $H$ base-policy actions, a direct decision evaluation is $O(HB^{d+1})$ in the worst case; memoization is effective when different prefixes reach the same mask state. The reported $d=3$ setting is consequently an offline planner. Runtime is measured, but it is not used as evidence of real-time capability.

\section{Experimental Protocol}
\label{sec:simulation}

\subsection{Frozen ISS-Mesh Decision Problem}

The benchmark uses the public NASA ISS glTF mesh \cite{nasa_iss_gltf_model}. The transformed mesh contains 247,525 triangles and is sampled into 90 equal-represented-area targets. Candidate generation at a nominal \unit{34}{m} radius with stride two yields 41 targets visible from at least one retained observation. A deterministic coverage-preserving reduction freezes 24 observation nodes. The archive contains 24 initial-to-observation records and $24\times23$ inter-observation records, for 576 directed edges in total. Every method receives this same archive. Each rest-to-rest segment lasts \unit{90}{s} and uses fixed-step RK4 propagation with $\dt=\unit{3}{s}$. The HCW mean motion is $1.13137\times10^{-3}$~rad/s, $u_{\max}=0.060~\mathrm{m/s^2}$, $v_{\max}=\unit{2.0}{m/s}$, $d_{\mathrm{safe}}=\unit{2}{m}$, $\varepsilon_r=\unit{0.5}{m}$, and $\varepsilon_v=\unit{0.05}{m/s}$. Equation~\eqref{eq:edge-cost} uses $(\lambda_v,\lambda_e,\lambda_s,\lambda_a)=(5,0.18,120,0.05)$.

Every method stops when it first reaches $C^{\mathrm{insp}}\ge0.80$ and may select at most 14 SOOAs. Feasibility requires every selected edge to satisfy the input, speed, full-mesh sampled-clearance, swept mesh-nonintersection, and terminal checks. Passive drift is disabled ($T_p=0$), so passive safety is not inferred. The physical $\Delta v$, minimum clearance above the \unit{2}{m} margin, and peak input are accumulated from the same archived edge records used by the graph objective.

\subsection{Split Discipline and Structure Selection}

The protocol was frozen before test evaluation. Deterministic, disjoint seed ranges define 24 training, 12 validation, 30 test, and 20 distribution-shift scenarios. Training scenarios may fit critic parameters only; validation scenarios select the algorithmic structure and rollout depth; test scenarios support the primary claim; shifted scenarios are inspected only afterward. A scenario removes 0--10\% of graph nodes and multiplies the weights of 6--14\% of targets by a factor sampled from $[2,4]$. The shifted split increases node removal to 15--25\%, changes 12--20\% of target priorities, and uses multipliers in $[5,8]$. Scenarios for which the deterministic incumbent cannot reach the goal are rejected during generation. The resulting inference is therefore conditional on incumbent feasibility, an explicit limitation rather than a population success-rate estimate.

Development considered a frozen fitted critic, a hand-scaled critic heuristic, a fixed-incumbent-suffix rollout, and repeated adaptive rollout. On the full 12-scenario validation set, adaptive depths one, two, and three all reached 12/12 goals; their mean graph costs were 101.424, 97.388, and 93.586, compared with 109.299 for local search. The predeclared lexicographic rule prioritized success, paired penalized cost, then latency. Depth three won all 12 cost pairs and was frozen before either test split was opened. Figure~\ref{fig:adp-development} reports this validation-only selection.

\begin{figure*}[t]
  \centering
  \includegraphics[width=\textwidth]{figures/adp_future/adp_development_selection.pdf}
  \caption{Validation-only structure selection. All adaptive rollout depths retained mission success; increasing the exact prefix from one to three reduced validation graph cost at increased online computation. Depth three was frozen before held-out evaluation. Validation outcomes are development evidence, not confirmatory test evidence.}
  \label{fig:adp-development}
\end{figure*}

Table~\ref{tab:baseline-definitions} separates the proposed adaptive rollout from the local, fixed-suffix, incumbent, fitted-value, and heuristic alternatives used in the paired study.

\begin{table*}[t]
  \centering
  \caption{Paired methods evaluated under identical geometry, dynamics, shield, and stopping rules.}
  \label{tab:baseline-definitions}
  \renewcommand{\arraystretch}{1.08}
  \setlength{\tabcolsep}{5pt}
  \footnotesize
  \begin{tabularx}{\textwidth}{@{}p{0.20\textwidth}p{0.34\textwidth}X@{}}
    \toprule
    Method & Selection rule & Purpose in the comparison \\
    \midrule
    Proposed rollout ADP & Algorithm~\ref{alg:safe-graph-api} with adaptive base completion and frozen depth $d=3$; no learned critic or incumbent fallback. & Tests the main ADP contribution. \\
    Local search & Three audited best-improvement reversal passes from the incumbent. & Strongest non-ADP sequence baseline and primary comparator. \\
    Fixed rollout & One-step rollout using a fixed incumbent suffix, with incumbent safeguard. & Separates adaptive recomputation from conventional fixed-suffix rollout. \\
    Incumbent & Deterministic shielded gain-efficiency construction without reversal passes. & Feasible reference policy. \\
    Frozen fitted critic & Ridge-fitted 26-feature value approximation with depth-two lookahead and no safeguard. & Tests whether function approximation alone is adequate. \\
    Heuristic search & Depth-two search using the hand-scaled initial value features and no safeguard. & Separates lookahead from fitted values. \\
    \bottomrule
  \end{tabularx}
\end{table*}

\subsection{Predeclared Held-Out Inference}

The primary endpoint is the paired penalized graph-cost difference $D_J=J_{\mathrm{ADP}}-J_{\mathrm{local}}$ on the 30 test scenarios. Penalized cost adds a fixed coverage-shortfall and mission-failure penalty, preventing a short unsuccessful route from appearing superior; because both primary methods succeed in all test pairs, it equals audited graph cost in the reported comparison. The claim gate requires no lower success rate than local search and an upper 95\% paired-bootstrap confidence bound below zero. Physical $\Delta v$ is the principal secondary endpoint. We report paired mean and median differences, 10,000-resample percentile intervals for the mean, Cohen's $d_z$, wins/ties/losses, and an exact two-sided sign test after excluding numerical ties. Secondary baseline tests use Holm correction, with adjusted results retained in the statistical archive. Coverage, action count, safety audits, and online time are reported descriptively; time is environment-sensitive and not a success gate.

\subsection{Saved Results and Reproducibility Boundary}

The result bundle contains the configuration snapshot, graph and scenario archives, one row per method--scenario pair, statistical output, figure source tables, and script hashes. A second run loads the frozen graph and repeats all 350 method evaluations. Reproducibility is declared only if graph content, scenarios, routes, success flags, coverage, costs, and physical audit metrics match exactly; wall-clock fields are excluded. Figure captions are linked to source files and transformations in a machine-readable trace manifest. The case-study generator reconstructs its routes from the archived scenario and aborts if route cost differs from the held-out row.

\subsection{ROS 2 Execution Interface and Deferred Validation}

The offline records connect to a ROS~2 Jazzy execution stack without making Gazebo physics the source of truth. The trajectory replay node reads the saved state, viewpoint, and attitude streams and publishes position, state, and attitude references on dedicated chaser topics. A tracking controller produces a nominal acceleration, the safety filter publishes the command consumed by the HCW dynamics node, and the dynamics node publishes chaser odometry. The evaluation node records trajectory, control, coverage, safety, planner, and mission-event CSV files together with a rosbag; RViz2 and Gazebo Harmonic visualize the same ROS state. This interface preserves the intended separation between planning, closed-loop execution, state propagation, logging, and visualization.

The baseline Jazzy workspace previously completed a short infrastructure smoke run, but that run predates the current SOOA graph planner and is not presented as validation. Closed-loop replay of the frozen routes is deferred until the target Ubuntu~24.04 and ROS~2 Jazzy environment is available. The replay protocol will report planned-versus-executed position error, realized $\Delta v$, minimum executed clearance, safety-filter intervention count, achieved coverage, dropped-reference count, and completion status. These quantities lie outside the present evidence boundary; no synthetic ROS measurements are introduced into the offline comparisons.

\section{Results}
\label{sec:results}

\subsection{Held-Out Mission Performance}

The frozen depth-3 rollout policy reached the 80\% goal in all 30 test scenarios and all 20 shifted scenarios. The local, fixed-rollout, and incumbent baselines also had complete success, making their raw graph costs directly comparable. Table~\ref{tab:paired-aggregate} shows that rollout ADP has the lowest mean graph cost and $\Delta v$ on both splits. On test, it selected the same mean number of actions as local search (10.8) but used a different sequence; its mean terminal coverage is slightly lower because it avoids surplus coverage after passing the common threshold. Figure~\ref{fig:paired-performance} shows the paired all-scenario comparisons rather than only aggregate means.

\begin{table*}[t]
  \centering
  \caption{Aggregate held-out results. Every listed method reaches the goal in every scenario. Coverage is weighted inspectable coverage at termination.}
  \label{tab:paired-aggregate}
  \renewcommand{\arraystretch}{1.06}
  \setlength{\tabcolsep}{5pt}
  \footnotesize
  \begin{tabular}{@{}llrrrrr@{}}
    \toprule
    Split & Method & Success & $C^{\mathrm{insp}}$ (\%) & Mean SOOAs & Mean $\Delta v$ (m/s) & Mean $J$ \\
    \midrule
    Test & Proposed rollout ADP & 30/30 & 81.00 & 10.80 & 11.239 & 92.789 \\
         & Local search & 30/30 & 81.39 & 10.80 & 12.946 & 106.804 \\
         & Fixed rollout & 30/30 & 81.26 & 10.73 & 12.518 & 103.286 \\
         & Incumbent & 30/30 & 81.67 & 11.30 & 13.566 & 111.924 \\
    \addlinespace
    Shifted & Proposed rollout ADP & 20/20 & 80.61 & 9.00 & 11.162 & 92.063 \\
            & Local search & 20/20 & 81.21 & 9.15 & 12.182 & 100.453 \\
            & Fixed rollout & 20/20 & 80.97 & 9.15 & 11.979 & 98.782 \\
            & Incumbent & 20/20 & 81.28 & 9.40 & 13.083 & 107.853 \\
    \bottomrule
  \end{tabular}
\end{table*}

\begin{figure*}[!t]
  \centering
  \includegraphics[width=\textwidth]{figures/adp_future/adp_heldout_performance.pdf}
  \caption{Held-out paired performance of proposed rollout ADP versus audited local search. (a) Test $\Delta v$, (b) test graph cost, and (c)--(d) the corresponding shifted-scenario comparisons. Dashed lines denote parity. Rollout ADP wins 28 of 30 test pairs with two numerical ties and wins 17 of 20 shifted pairs with one tie. Insets report paired-bootstrap intervals for the mean difference; coverage success is matched in every panel.}
  \label{fig:paired-performance}
\end{figure*}

\subsection{Paired Effect Estimates}

The primary claim gate is passed. On the test split, the ADP-minus-local graph-cost difference is $-14.016$ with 95\% interval $[-18.265,-10.451]$ and $d_z=-1.274$. This is a 13.12\% reduction relative to the local mean. The corresponding physical $\Delta v$ difference is \unit{-1.707}{m/s}, interval $[-2.199,-1.265]$, $d_z=-1.274$, or 13.18\%. Both endpoints have 28 wins, two numerical ties, no losses, and exact sign-test $p=7.45\times10^{-9}$. Mean selected-count difference is zero. Thus the advantage is route quality, not a comparison between a short failed path and a complete mission. Table~\ref{tab:paired-inference} gives the paired estimates for both evaluation splits.

The shifted split is directionally consistent but is treated as a stress test rather than a second confirmatory sample. Graph cost falls by 8.35\% and $\Delta v$ by 8.37\%; the latter paired mean is \unit{-1.020}{m/s}, with interval $[-1.417,-0.635]$, 17 wins, one tie, two losses, and $p=7.29\times10^{-4}$. The effect is smaller under the stronger target-priority and node-dropout shift, which is more informative than claiming uniform dominance.

\begin{table*}[t]
  \centering
  \caption{Paired rollout-ADP minus local-search inference. Confidence intervals use 10,000 paired bootstrap resamples. The exact sign test excludes numerical ties.}
  \label{tab:paired-inference}
  \renewcommand{\arraystretch}{1.06}
  \setlength{\tabcolsep}{5pt}
  \footnotesize
  \begin{tabular}{@{}lrrrrr@{}}
    \toprule
    Split and endpoint & Mean difference & Median difference & 95\% CI for mean & Win/tie/loss & Sign-test $p$ \\
    \midrule
    Test: graph cost $J$ & $-14.016$ & $-12.842$ & $[-18.265,-10.451]$ & 28/2/0 & $7.45\times10^{-9}$ \\
    Test: $\Delta v$ (m/s) & $-1.707$ & $-1.556$ & $[-2.199,-1.265]$ & 28/2/0 & $7.45\times10^{-9}$ \\
    Shifted: graph cost $J$ & $-8.390$ & $-7.215$ & $[-11.677,-5.223]$ & 17/1/2 & $7.29\times10^{-4}$ \\
    Shifted: $\Delta v$ (m/s) & $-1.020$ & $-0.873$ & $[-1.417,-0.635]$ & 17/1/2 & $7.29\times10^{-4}$ \\
    \bottomrule
  \end{tabular}
\end{table*}

\subsection{Component and Safety Audit}

The negative critic result clarifies the contribution. The frozen fitted critic succeeds in only 18/30 test and 11/20 shifted scenarios; the unfitted heuristic reaches 25/30 and 13/20. In contrast, fixed rollout, local search, and adaptive rollout all achieve complete success. The learned value approximation is therefore neither concealed nor used to label the successful algorithm ``ADP.'' The successful route is rollout ADP because it uses a base-policy cost-to-go approximation inside repeated truncated Bellman improvement.

All selected edges share the same shield. Across the proposed test and shifted routes, the worst minimum clearance is \unit{0.0619}{m} above the required \unit{2}{m} mesh distance. Peak proposed inputs are 0.0553 and $0.0556~\mathrm{m/s^2}$ on the two splits, below the $0.060~\mathrm{m/s^2}$ limit. These are sampled deterministic audits; passive drift is disabled and continuous-time robustness is not claimed. Median test online time is \unit{1.64}{s} for rollout ADP versus \unit{0.0076}{s} for local search on the development computer, exposing the expected depth-three computation cost. Figure~\ref{fig:paired-safety} juxtaposes this safety evidence with the unsuccessful critic ablations.

\begin{figure*}[!t]
  \centering
  \includegraphics[width=\textwidth]{figures/adp_future/adp_ablation_safety.pdf}
  \caption{Component and safety audit. (a) Test $\Delta v$ distributions with bootstrap intervals for method means. (b) Mission success exposes the failure of standalone fitted and hand-scaled critics; fixed rollout, local search, and proposed adaptive rollout reach every goal. (c) Minimum sampled clearance above the \unit{2}{m} margin and (d) peak input for ADP versus local search. Red lines are acceptance limits. Passive drift and uncertainty are not evaluated.}
  \label{fig:paired-safety}
\end{figure*}

\subsection{Representative Generated-Trajectory Case}

To show the route-level mechanism, the figure generator selects test scenario~2 post hoc as the paired graph-cost difference closest to the test median. This outcome-based rule makes the case representative of the observed effect size, but also makes it illustrative rather than inferential. The archived scenario seed is 20733, and reconstructed route costs match the held-out rows exactly.

Rollout ADP reaches 81.79\% coverage with 10 SOOAs, graph cost 99.204, and \unit{12.026}{m/s} $\Delta v$. Local search reaches 83.82\% with 11 SOOAs, cost 111.966, and \unit{13.574}{m/s} $\Delta v$. Figure~\ref{fig:representative-trajectory} shows that ADP does not merely delete the last local action. It selects a different early ordering, accumulates useful coverage faster, and crosses the 80\% goal one action earlier. The local route's surplus terminal coverage does not compensate for its larger cost under the common thresholded objective.

\begin{figure*}[!t]
  \centering
  \includegraphics[width=\textwidth]{figures/adp_future/adp_representative_trajectory.pdf}
  \caption{Illustrative median-effect held-out case. (a) Generated observation-to-observation routes in the LVLH frame; markers are terminal observations and the star is the common initial state. (b) Weighted inspectable coverage versus executed SOOA count. (c) Cumulative physical $\Delta v$. The ADP route reaches the goal after 10 actions; local search uses 11. The selection is post hoc and does not replace the all-scenario inference.}
  \label{fig:representative-trajectory}
\end{figure*}

\subsection{Reproducibility and Evidence Boundary}

The frozen-graph rerun reproduces every non-timing field in all 350 method--scenario rows, including selected node sequences, success, coverage, graph cost, $\Delta v$, clearance, and peak input. The result supports a bounded statement: within the frozen SOOA library and incumbent-feasible scenario generator, depth-three rollout ADP consistently improves the strongest local baseline while retaining sampled graph viability. It does not establish superiority for scenarios without a feasible base completion, for different meshes or transfer durations, or under navigation, actuator, and model uncertainty.

\section{Discussion}
\label{sec:discussion}

The central result is algorithmic rather than terminological. A fitted critic is not necessary for approximate dynamic programming: in rollout ADP, the complete cost of a tractable base policy is the approximate value, and Bellman minimization improves its early decisions \cite{bertsekas2020rollout}. This form is more defensible here than the fitted alternatives because the mission has a hard terminal coverage condition. An unconstrained critic may underestimate an incomplete path, whereas the proposed terminal evaluator assigns infinite value unless a shielded completion actually reaches the goal. The fitted critic's 60\% test success rate confirms that this is a material distinction, not only a theoretical preference.

The advantage over local sequence improvement is reach. Reversal search explores permutations near one incumbent and cannot introduce an observation absent from that sequence. Rollout ADP evaluates every currently safe observation, anticipates three exact actions, and then asks whether an adaptive base policy can complete the mission. The representative case shows the consequence: a different early sequence reaches the goal with one fewer action and lower $\Delta v$. Across the complete test set, mean action count is unchanged, yet $\Delta v$ falls by 13.18\%, showing that the dominant aggregate mechanism is better transition ordering rather than systematic route truncation.

The benefit has a computational price. The proposed median test planning time is approximately 216 times that of local search in the present Python implementation. This comparison is a diagnostic, not a real-time benchmark: graph records are cached, timing fields vary across reruns, and no optimized parallel expansion is used. Batched action evaluation, transposition tables retained across decisions, dominance pruning for coverage masks, and a learned value used to order (but not eliminate) viable branches are promising routes to reduce latency while preserving the finite-completion safeguard.

The paired result is stronger than a single favorable route but narrower than a general robustness claim. Scenario pairing varies target priorities and graph-node availability while holding geometry and edge dynamics fixed. The test and shifted confidence intervals quantify consistency over those designed changes, not a probability of flight improvement. Moreover, scenario generation accepts only instances for which the incumbent is successful. Proposition~\ref{prop:rollout-viability} therefore explains preservation of an available completion; it does not establish recovery when no feasible base completion exists. Such cases require a repair policy, a broader candidate library, or exact feasibility search.

The relation to formation-based observation is similarly complementary. Near-natural relative ellipses and closed-loop formation control around a rotating station can provide lower-effort motion structure and stronger stability analysis \cite{zhou2026formation}. OrbInspect instead resolves which surface regions are visible, how prescribed observations are credited, and which observation should follow the current state. A useful synthesis would generate SOOA candidates from near-natural orbit families, then select among them using the present coverage and camera model. This could reduce transfer effort while preserving explicit target-level inspection accounting.

The safety evidence remains limited to the stated model. HCW translation assumes a circular chief and a fixed LVLH target. Clearance samples and swept surface-intersection tests use the complete deterministic ISS mesh, but they do not create a closed-loop proof under navigation error, target rotation, delayed commands, actuator uncertainty, or unmodeled disturbances. The study disables the optional passive-drift audit, so no passive-safety conclusion should be drawn. The full solver also has no continuous-space optimality guarantee.

The camera evidence has comparable boundaries. The visibility model includes range, field of view, incidence, and mesh occlusion, but the offline graph credits a prescribed terminal observation rather than simulating exposure interruption or closed-loop dwell. Illumination, blur, plume contamination, camera slew, and detailed attitude dynamics are omitted. ROS~2 and Gazebo replay remain implementation work and are not treated as validation in the present manuscript.

A flight-oriented development path should first validate closed-loop replay on Ubuntu~24.04 with ROS~2 Jazzy. The next algorithmic study should deliberately include incumbent-infeasible scenarios and compare viability repair, beam rollout, and constrained fitted values under a fixed test protocol. Batched full-mesh edge evaluation and cached passive audits can then support passive-drift experiments. Uncertainty-inflated keep-out volumes, camera-slew and illumination constraints, navigation error, a rotating body frame, and an optional Basilisk backend should follow while retaining HCW as a transparent regression reference.

\section{Conclusion}
\label{sec:conclusion}

This paper formulated orbital inspection as viability-preserving rollout approximate dynamic programming over finite-duration observable orbital arcs. The proposed policy retains the current node, covered and selected masks, and remaining budget; expands a finite exact Bellman prefix; and uses a complete adaptive safe-greedy policy as its terminal value. Infinite cost for unsuccessful completions prevents approximate values from trading away the mission goal. Full-mesh queries supply visibility, sampled clearance, and swept surface-crossing checks from one consistent ISS representation.

On 30 frozen, incumbent-feasible test scenarios, rollout ADP and local search both reached the 80\% inspectable-coverage goal. ADP reduced mean $\Delta v$ by 13.18\%, from 12.946 to \unit{11.239}{m/s}; the paired mean difference was \unit{-1.707}{m/s}, with a 95\% interval $[-2.199,-1.265]$, 28 wins, two ties, and no losses. An 8.37\% reduction persisted under stronger node dropout and target-priority shift. The fitted-critic route failed the success gate, while the model-based rollout structure was exactly reproducible in all non-timing fields. The evidence therefore supports rollout ADP as the main contribution, rather than the fitted critic or local reversal search. The conclusion remains limited to a deterministic sampled graph and does not claim passive safety, continuous-time robustness, continuous-space optimality, or ROS execution performance.

\section*{Acknowledgment}
This research was supported in part by the Research Centre for Unmanned Autonomous Systems of The Hong Kong Polytechnic University under Grant P0046487. OpenAI Codex was used during manuscript preparation to assist with language editing across all sections, LaTeX formatting, source-checking support, code review, and development of figure-generation scripts. The authors reviewed and verified the mathematical statements, references, source code, figures, and reported results and take responsibility for the content.

\section*{Data and Code Availability}
The complete reproducibility package accompanies this submission as supplementary material. It contains the experiment configuration, source code, frozen graph and scenarios, raw method rows, statistical outputs, figure source tables, transformation hashes, and deterministic rerun report. A mirrored copy is available from the corresponding author upon reasonable request.

\bibliographystyle{IEEEtran}
\bibliography{references}

\begin{IEEEbiography}
    [{\includegraphics[width=1in,height=1.25in,clip,keepaspectratio]{figures/authors/trg.jpg}}]{Rugang Tang} received the B.Eng. degree in aircraft design and engineering from Northwestern Polytechnical University, Xi'an, China, in 2020. He is pursuing the Ph.D. degree in aircraft design with the Department of Aeronautical and Aviation Engineering, The Hong Kong Polytechnic University, Hong Kong. His research interests include multi-agent and nonlinear systems, reach-avoid differential graphical games, and approximate dynamic programming.
\end{IEEEbiography}

\begin{IEEEbiography}[{
\includegraphics[width=1in,height=1.25in,clip,keepaspectratio]{figures/authors/nx.pdf}}]{Xin Ning} received the B.Eng., M.Eng., and Ph.D. degrees from Northwestern Polytechnical University, Xi'an, China, in 2004, 2007, and 2011, respectively. He is a Professor and Vice Dean of the School of Astronautics at Northwestern Polytechnical University and was a visiting scholar at the University of Calgary. His research interests include aerospace flight dynamics, navigation, guidance and control, spacecraft on-orbit servicing, trajectory planning, intelligent optimization, nonlinear systems, and modeling.
\end{IEEEbiography}

\begin{IEEEbiography}
    [{
\includegraphics[width=1in,height=1.25in,clip,keepaspectratio]{figures/authors/CYW.png}}]{Chih-Yung Wen} received the B.S. degree from National Taiwan University in 1986 and the M.S. and Ph.D. degrees from the California Institute of Technology in 1989 and 1994, respectively. He is Head and Chair Professor of Aeronautical Engineering with the Department of Aeronautical and Aviation Engineering, The Hong Kong Polytechnic University, and Associate Director of its Research Institute for Sports Science and Technology. His professional activities span mechanical and aerospace engineering.
\end{IEEEbiography}

\end{document}
